\chapter{2008年大纲II卷（文）}

\section{单选题}

\begin{problem}
若 \(\sin \alpha < 0\) 且 \(\tan \alpha > 0\)，则 \(\alpha\) 是（\quad）

\choices
    {第一象限角}
    {第二象限角}
    {第三象限角}
    {第四象限角}
\end{problem}

\begin{answer}
C
\end{answer}

\begin{solution}
\(\sin\alpha<0\) 时，\(\alpha\) 位于第三或第四象限；\(\tan\alpha>0\) 时，\(\alpha\) 位于第一或第三象限.两个条件同时成立时，\(\alpha\) 只能位于第三象限，故选 C.
\end{solution}

\begin{problem}
设集合 \(M = \{m \in \Z \mid -3 < m < 2\}\)，\(N = \{n \in \Z \mid -1 \leqslant n \leqslant 3\}\)，则 \(M \cap N =\)（\quad）

\choices
    {\(\{0,1\}\)}
    {\(\{-1,0,1\}\)}
    {\(\{0,1,2\}\)}
    {\(\{-1,0,1,2\}\)}
\end{problem}

\begin{answer}
B
\end{answer}

\begin{solution}
因为 \(m\) 为整数且 \(-3<m<2\)，所以 \(M=\{-2,-1,0,1\}\).同理，\(N=\{-1,0,1,2,3\}\).取两集合的公共元素，得 \(M\cap N=\{-1,0,1\}\)，故选 B.
\end{solution}

\begin{problem}
原点到直线 \(x + 2y - 5 = 0\) 的距离为（\quad）

\choices
    {\(1\)}
    {\(\sqrt{3}\)}
    {\(2\)}
    {\(\sqrt{5}\)}
\end{problem}

\begin{answer}
D
\end{answer}

\begin{solution}
点 \((x_0,y_0)\) 到直线 \(Ax+By+C=0\) 的距离为 \(\dfrac{|Ax_0+By_0+C|}{\sqrt{A^2+B^2}}\).取原点 \((0,0)\)，得到距离
\(\dfrac{|0+0-5|}{\sqrt{1^2+2^2}}=\dfrac{5}{\sqrt{5}}=\sqrt{5}\)，故选 D.
\end{solution}

\begin{problem}
函数 \( f(x) = \frac{1}{x} - x \) 的图象关于（\quad）

\choices
    {\(y\) 轴对称}
    {直线 \(y = -x\) 对称}
    {坐标原点对称}
    {直线 \( y = x \) 对称}
\end{problem}

\begin{answer}
C
\end{answer}

\begin{solution}
函数定义域为 \(\{x\in\R\mid x\ne0\}\)，关于原点对称.对任意 \(x\ne0\)，有
\(f(-x)=-\dfrac{1}{x}+x=-f(x)\)，所以 \(f(x)\) 是奇函数，其图象关于坐标原点对称，故选 C.
\end{solution}

\begin{problem}
若 \(x \in (\e^{-1}, 1)\)，\(a = \ln x\)，\(b = 2 \ln x\)，\(c = \ln^3 x\)，则（\quad）

\choices
    {\(a < b < c\)}
    {\(c < a < b\)}
    {\(b < a < c\)}
    {\(b < c < a\)}
\end{problem}

\begin{answer}
C
\end{answer}

\begin{solution}
令 \(t=\ln x\).由于 \(\e^{-1}<x<1\)，且 \(\ln x\) 在正数上单调递增，所以 \(-1<t<0\).于是 \(b-a=t<0\)，从而 \(b<a\).又
\(a-c=t-t^3=t(1-t^2)<0\)，因为 \(t<0\) 且 \(1-t^2>0\)，所以 \(a<c\).因此 \(b<a<c\)，故选 C.
\end{solution}

\begin{problem}
设变量 \(x\)，\(y\) 满足约束条件：\(\begin{cases}
y \geqslant x,\\
x + 2y \leqslant 2,\\
x \geqslant -2
\end{cases}\)，则 \(z = x - 3y\) 的最小值（\quad）

\choices
    {\(-2\)}
    {\(-4\)}
    {\(-6\)}
    {\(-8\)}
\end{problem}

\begin{answer}
D
\end{answer}

\begin{solution}
由 \(y\geqslant x\) 和 \(x+2y\leqslant2\)，得 \(x\leqslant\dfrac{2-x}{2}\)，即 \(x\leqslant\dfrac23\).结合 \(x\geqslant-2\)，可取 \(-2\leqslant x\leqslant\dfrac23\).

固定 \(x\) 时，\(z=x-3y\) 随 \(y\) 增大而减小，所以在约束范围内取 \(y=\dfrac{2-x}{2}\).此时
\(z=x-3\cdot\dfrac{2-x}{2}=\dfrac{5x-6}{2}\)，它随 \(x\) 增大而增大，故在 \(x=-2\) 时取得最小值.相应地 \(y=2\)，所以 \(z_{\min}=\dfrac{5(-2)-6}{2}=-8\)，故选 D.
\end{solution}

\begin{problem}
设曲线 \(y = ax^2\) 在点 \((1, a)\) 处的切线与直线 \(2x - y - 6 = 0\) 平行，则 \(a =\)（\quad）

\choices
    {\(1\)}
    {\(\frac{1}{2}\)}
    {\(-\frac{1}{2}\)}
    {\(-1\)}
\end{problem}

\begin{answer}
A
\end{answer}

\begin{solution}
曲线的导数为 \(y'=2ax\)，所以在 \(x=1\) 处的切线斜率为 \(2a\).直线 \(2x-y-6=0\) 化为 \(y=2x-6\)，斜率为 \(2\).由两直线平行知 \(2a=2\)，解得 \(a=1\)，故选 A.
\end{solution}

\begin{problem}
正四棱锥的侧棱长为 \(2\sqrt{3}\)，侧棱与底面所成的角为 \(60^{\circ}\)，则该棱锥的体积为（\quad）

\choices
    {\(3\)}
    {\(6\)}
    {\(9\)}
    {\(18\)}
\end{problem}

\begin{answer}
B
\end{answer}

\begin{solution}
设棱锥高为 \(h\)，顶点在正方形底面的射影为底面中心 \(O\)，底面一条顶点到中心的距离为 \(r\).侧棱长为 \(2\sqrt3\)，侧棱与底面成角 \(60^{\circ}\)，故
\(h=2\sqrt3\sin60^{\circ}=3\)，\(r=2\sqrt3\cos60^{\circ}=\sqrt3\).

正方形底面边长设为 \(s\)，则其中心到顶点的距离为 \(\dfrac{s\sqrt2}{2}\)，所以 \(\dfrac{s\sqrt2}{2}=\sqrt3\)，得 \(s^2=6\).底面积为 \(6\)，从而棱锥体积为 \(\dfrac13\times6\times3=6\)，故选 B.
\end{solution}

\begin{problem}
\((1 - \sqrt{x})^4 (1 + \sqrt{x})^4\) 的展开式中 \(x\) 的系数是（\quad）

\choices
    {\(-4\)}
    {\(-3\)}
    {\(3\)}
    {\(4\)}
\end{problem}

\begin{answer}
A
\end{answer}

\begin{solution}
利用平方差公式，原式为 \(\bigl((1-\sqrt{x})(1+\sqrt{x})\bigr)^4=(1-x)^4\).展开得
\(1-4x+6x^2-4x^3+x^4\)，所以 \(x\) 的系数为 \(-4\)，故选 A.
\end{solution}

\begin{problem}
函数 \( f(x) = \sin x - \cos x \) 的最大值为（\quad）

\choices
    {\(1\)}
    {\(\sqrt{2}\)}
    {\(\sqrt{3}\)}
    {\(2\)}
\end{problem}

\begin{answer}
B
\end{answer}

\begin{solution}
由辅助角公式，\(\sin x-\cos x=\sqrt2\sin\left(x-\dfrac\pi4\right)\).正弦函数的最大值为 \(1\)，所以原函数的最大值为 \(\sqrt2\)，故选 B.
\end{solution}

\begin{problem}
设 \(\triangle ABC\) 是等腰三角形，\(\angle ABC = 120^{\circ}\)，则以 \(A\)，\(B\) 为焦点且过点 \(C\) 的双曲线的离心率为（\quad）

\choices
    {\(\frac{1 + \sqrt{2}}{2}\)}
    {\(\frac{1 + \sqrt{3}}{2}\)}
    {\(1 + \sqrt{2}\)}
    {\(1 + \sqrt{3}\)}
\end{problem}

\begin{answer}
B
\end{answer}

\begin{solution}
由于三角形的一个角为 \(120^{\circ}\)，它不可能是等腰三角形的底角，因此 \(AB=BC\).设 \(AB=BC=l\)，由余弦定理得
\(AC^2=l^2+l^2-2l^2\cos120^{\circ}=3l^2\)，所以 \(AC=\sqrt3l\).

双曲线的焦点为 \(A\)、\(B\)，且过 \(C\)，故焦距之差的绝对值为 \(AC-BC=(\sqrt3-1)l\).若双曲线的实半轴为 \(a_0\)，焦半距为 \(c_0\)，则 \(2a_0=(\sqrt3-1)l\)，\(2c_0=AB=l\).因此离心率
\(e=\dfrac{c_0}{a_0}=\dfrac{1}{\sqrt3-1}=\dfrac{1+\sqrt3}{2}\)，故选 B.
\end{solution}

\begin{problem}
已知球的半径为 \(2\)，相互垂直的两个平面分别截球面而得两个圆. 若两圆的公共弦长为 \(2\)，则两圆的圆心距等于（\quad）

\choices
    {\(1\)}
    {\(\sqrt{2}\)}
    {\(\sqrt{3}\)}
    {\(2\)}
\end{problem}

\begin{answer}
C
\end{answer}

\begin{solution}
设两平面的交线为 \(z\) 轴，两个平面分别为 \(x=0\) 和 \(y=0\)，球心记为 \(O=(u,v,w)\).两个截面圆的圆心分别是 \(O_1=(0,v,w)\) 和 \(O_2=(u,0,w)\)，所以
\(O_1O_2^2=u^2+v^2\).两平面的交线与球面的交点是公共弦的端点，公共弦的中点为 \((0,0,w)\).由球半径为 \(2\)、公共弦长为 \(2\)，公共弦半长为 \(1\)，故
\(u^2+v^2=2^2-1^2=3\).于是 \(O_1O_2=\sqrt3\)，故选 C.
\end{solution}

\section{填空题}

\begin{problem}
设向量 \(\bs{a} = (1, 2)\)，\(\bs{b} = (2, 3)\). 若向量 \(\lambda\bs{a} + \bs{b}\) 与向量 \(\bs{c} = (-4, -7)\) 共线，则 \(\lambda =\fillinblank{}\).
\end{problem}

\begin{answer}
\(2\)
\end{answer}

\begin{solution}
\(\lambda\bs{a}+\bs{b}=(\lambda+2,2\lambda+3)\).两个向量共线的充要条件是坐标行列式为零，所以
\((\lambda+2)(-7)-(2\lambda+3)(-4)=0\).化简得 \(\lambda-2=0\)，从而 \(\lambda=2\).
\end{solution}

\begin{problem}
从 \(10\text{ 名}\) 男同学，\(6\text{ 名}\) 女同学中选 \(3\text{ 名}\) 参加体能测试，则选到的 \(3\text{ 名}\) 同学中既有男同学又有女同学的不同选法共有\fillinblank{} 种.（用数字作答）
\end{problem}

\begin{answer}
\(420\)
\end{answer}

\begin{solution}
选出的三人中既有男生又有女生，分为选一名男生和两名女生，或选两名男生和一名女生两类.选法总数为
\(\mathrm C_{10}^{1}\mathrm C_{6}^{2}+\mathrm C_{10}^{2}\mathrm C_{6}^{1}=10\times15+45\times6=150+270=420\)，所以应填 \(420\).
\end{solution}

\begin{problem}
已知 \(F\) 是抛物线 \(C: y^{2} = 4x\) 的焦点，\(A\)，\(B\) 是 \(C\) 上的两个点，线段 \(AB\) 的中点为 \(M(2,2)\)，则 \(\triangle ABF\) 的面积等于\fillinblank{}.
\end{problem}

\begin{answer}
\(2\)
\end{answer}

\begin{solution}
设 \(A=(x_1,y_1)\)，\(B=(x_2,y_2)\).由中点为 \(M(2,2)\)，得 \(x_1+x_2=4\)，\(y_1+y_2=4\).又因 \(A\)、\(B\) 在 \(y^2=4x\) 上，作差得
\((y_1-y_2)(y_1+y_2)=4(x_1-x_2)\).代入 \(y_1+y_2=4\)，得到 \(y_1-y_2=x_1-x_2\)，故弦 \(AB\) 的斜率为 \(1\).

直线 \(AB\) 经过 \((2,2)\)，所以其方程为 \(y=x\).与抛物线联立得 \(x^2=4x\)，从而交点为 \(A=(0,0)\)、\(B=(4,4)\).抛物线 \(y^2=4x\) 的焦点为 \(F=(1,0)\)，于是
\(S_{\triangle ABF}=\dfrac12\left|\det\begin{pmatrix}4&4\\1&0\end{pmatrix}\right|=2\).
\end{solution}

\begin{problem}
平面内的一个四边形为平行四边形的充要条件有多个，如两组对边分别平行，类似地，写出空间中的一个四棱柱为平行六面体的两个充要条件：

\begin{circlelist}
    \item\fillinblank{}；
    \item\fillinblank{}.
\end{circlelist}

（写出你认为正确的两个充要条件）
\end{problem}

\begin{answer}
三组相对侧面分别平行；四条体对角线交于一点且在交点处互相平分.
\end{answer}

\begin{solution}
以下给出两组充要条件.第一组是“三组相对侧面分别平行”.平行六面体的三组相对侧面分别平行；反过来，四棱柱的三组相对侧面分别平行时，其底面内的两组相对边分别平行，底面是平行四边形，因此该四棱柱是平行六面体.

第二组是“四条体对角线交于一点且在交点处互相平分”.平行六面体的四条体对角线具有共同的中点；反过来，四棱柱的四条体对角线有共同中点时，底面两条对角线互相平分，底面为平行四边形，因此该四棱柱是平行六面体.
\end{solution}

\section{解答题}

\begin{problem}
在 \(\triangle ABC\) 中，\(\cos A = -\frac{5}{13}\)，\(\cos B = \frac{3}{5}\).

\begin{enumerate}
\item 求 \(\sin C\) 的值；
\item 设 \(BC = 5\)，求 \(\triangle ABC\) 的面积.
\end{enumerate}
\end{problem}

\begin{solution}
\begin{enumerate}
\item 因为 \(A\)、\(B\) 是三角形内角，所以 \(\sin A>0\)，\(\sin B>0\).由已知得 \(\sin A=\sqrt{1-\cos^2A}=\dfrac{12}{13}\)，\(\sin B=\sqrt{1-\cos^2B}=\dfrac45\).又 \(C=\pi-(A+B)\)，所以
\(\sin C=\sin(A+B)=\sin A\cos B+\cos A\sin B=\dfrac{12}{13}\cdot\dfrac35-\dfrac5{13}\cdot\dfrac45=\dfrac{16}{65}\).
\item 由正弦定理，\(\dfrac{AB}{\sin C}=\dfrac{BC}{\sin A}\)，故 \(AB=5\cdot\dfrac{16}{65}\cdot\dfrac{13}{12}=\dfrac43\).因此
\(S_{\triangle ABC}=\dfrac12\,AB\cdot BC\cdot\sin B=\dfrac12\cdot\dfrac43\cdot5\cdot\dfrac45=\dfrac83\).
\end{enumerate}
\end{solution}

\begin{problem}
等差数列 \(\{a_{n}\}\) 中，\(a_{4} = 10\)，且 \(a_{3}\)，\(a_{6}\)，\(a_{10}\) 成等比数列. 求数列 \(\{a_{n}\}\) 前 \(20\) 项的和 \(S_{20}\).
\end{problem}

\begin{solution}
设等差数列的公差为 \(d\).由 \(a_4=10\)，有 \(a_3=10-d\)，\(a_6=10+2d\)，\(a_{10}=10+6d\).三项成等比数列的充要条件为中间项的平方等于两端项的积，因此
\((10+2d)^2=(10-d)(10+6d)\).展开化简得 \(10d^2-10d=0\)，所以 \(d=0\) 或 \(d=1\).

当 \(d=0\) 时，\(a_1=10\)，\(S_{20}=20\times10=200\).当 \(d=1\) 时，\(a_1=a_4-3d=7\)，故 \(S_{20}=\dfrac{20}{2}\bigl(2\times7+19\times1\bigr)=330\).所以 \(S_{20}=200\) 或 \(330\).
\end{solution}

\begin{problem}
甲，乙两人进行射击比赛. 在一轮比赛中，甲，乙各射击 \(1\text{ 发}\) 子弹. 根据以往资料知，甲击中 \(8\text{ 环}\)，\(9\text{ 环}\)，\(10\text{ 环}\) 的概率分别为 \(0.6\)，\(0.3\)，\(0.1\)，乙击中 \(8\text{ 环}\)，\(9\text{ 环}\)，\(10\text{ 环}\) 的概率分别为 \(0.4\)，\(0.4\)，\(0.2\). 设甲，乙的射击相互独立.

\begin{enumerate}
\item 求在一轮比赛中甲击中的环数多于乙击中环数的概率；
\item 求在独立的 \(3\) 轮比赛中，至少有 \(2\) 轮甲击中的环数多于乙击中环数的概率.
\end{enumerate}
\end{problem}

\begin{solution}
\begin{enumerate}
\item 设一轮比赛中甲击中的环数多于乙击中的环数为事件 \(A\).事件 \(A\) 包含三种互斥情形：甲击中 \(9\) 环而乙击中 \(8\) 环，甲击中 \(10\) 环而乙击中 \(8\) 环，甲击中 \(10\) 环而乙击中 \(9\) 环.由独立性，
\(P(A)=0.3\times0.4+0.1\times0.4+0.1\times0.4=0.2\).
\item 设三轮中甲击中环数多于乙击中环数的轮数为 \(X\).由各轮相互独立，\(X\) 服从参数为 \(3\)、\(0.2\) 的二项分布.因此
\(P(X\geqslant2)=\mathrm C_3^2(0.2)^2(0.8)+(0.2)^3=0.096+0.008=0.104\).
\end{enumerate}
\end{solution}

\begin{problem}
如图，正四棱柱 \(ABCD - A_{1}B_{1}C_{1}D_{1}\) 中，\(AA_{1} = 2AB = 4\)，点 \(E\) 在 \(CC_{1}\) 上且 \(C_{1}E = 3EC\).

\begin{enumerate}
\item 证明：\(A_{1}C \perp\) 平面 \(BED\)；
\item 求二面角 \(A_{1} - DE - B\) 的大小.
\end{enumerate}

\bitmapfigure[max width=0.34\linewidth]{img/2008/syllabus_paper_2_liberal/q20_fig1.png}
\end{problem}

\begin{solution}
建立空间直角坐标系，令
\(A=(0,0,0)\)，\(B=(2,0,0)\)，\(C=(2,2,0)\)，\(D=(0,2,0)\)，\(A_1=(0,0,4)\).由 \(C_1E=3EC\) 且 \(CC_1=4\)，得 \(EC=1\)，所以 \(E=(2,2,1)\).

\begin{enumerate}
\item 有 \(\overrightarrow{DB}=(2,-2,0)\)，\(\overrightarrow{DE}=(2,0,1)\)，\(\overrightarrow{A_1C}=(2,2,-4)\).于是
\(\overrightarrow{A_1C}\cdot\overrightarrow{DB}=2\times2+2\times(-2)+(-4)\times0=0\)，
\(\overrightarrow{A_1C}\cdot\overrightarrow{DE}=2\times2+2\times0+(-4)\times1=0\).直线 \(DB\)、\(DE\) 在平面 \(BED\) 内相交，且 \(A_1C\) 同时垂直于它们，所以 \(A_1C\perp\) 平面 \(BED\).
\item 取平面 \(A_1DE\) 和平面 \(BDE\) 的法向量
\(\bs{n}_1=\overrightarrow{DA_1}\times\overrightarrow{DE}=(-2,8,4)\)，\(\bs{n}_2=\overrightarrow{DB}\times\overrightarrow{DE}=(-2,-2,4)\).因此两平面的夹角 \(\theta\) 满足
\(\cos\theta=\dfrac{|\bs{n}_1\cdot\bs{n}_2|}{|\bs{n}_1|\,|\bs{n}_2|}=\dfrac{4}{\sqrt{84}\sqrt{24}}=\dfrac{\sqrt{14}}{42}\).取二面角的锐角，得 \(\theta=\arccos\dfrac{\sqrt{14}}{42}\).
\end{enumerate}
\end{solution}

\begin{problem}
设 \(\alpha \in \R\). 函数 \(f(x) = \alpha x^3 - 3x^2\).

\begin{enumerate}
\item 若 \(x = 2\) 是函数 \(y = f(x)\) 的极值点，求 \(\alpha\) 的值；
\item 若函数 \(g(x) = f(x) + f'(x)\)，\(x \in [0,2]\)，在 \(x = 0\) 处取得最大值，求 \(\alpha\) 的取值范围.
\end{enumerate}
\end{problem}

\begin{solution}
\begin{enumerate}
\item \(f'(x)=3\alpha x^2-6x\)，所以极值点 \(x=2\) 必须满足 \(f'(2)=12\alpha-12=0\)，得 \(\alpha=1\).此时 \(f'(x)=3x(x-2)\)，在 \(x=2\) 左侧为负、右侧为正，故 \(x=2\) 确为极小值点.
\item 由 \(f'(x)=3\alpha x^2-6x\)，有
\(g(x)=\alpha x^3+(3\alpha-3)x^2-6x=x\bigl(\alpha x^2+3(\alpha-1)x-6\bigr)\).又 \(g(0)=0\)，要使 \(x=0\) 在 \([0,2]\) 上取得最大值，必须且只需对所有 \(0<x\leqslant2\) 有 \(g(x)\leqslant0\)，即
\(\alpha\leqslant\varphi(x)=\dfrac{3x+6}{x^2+3x}=\dfrac{3(x+2)}{x(x+3)}\).

对 \(x>0\)，\(\varphi'(x)=\dfrac{3\bigl[x(x+3)-(x+2)(2x+3)\bigr]}{x^2(x+3)^2}<0\)，所以 \(\varphi(x)\) 在 \((0,2]\) 上递减，最小值为 \(\varphi(2)=\dfrac65\).因此条件等价于 \(\alpha\leqslant\dfrac65\).
\end{enumerate}
\end{solution}

\begin{problem}
设椭圆中心在坐标原点，\(A(2,0)\)，\(B(0,1)\) 是它的两个顶点，直线 \(y = kx\)（\(k > 0\)）与 \(AB\) 相交于点 \(D\)，与椭圆相交于 \(E\)，\(F\) 两点.

\begin{enumerate}
\item 若 \(\overrightarrow{ED} = 6\overrightarrow{DF}\)，求 \(k\) 的值；
\item 求四边形 \(AEBF\) 面积的最大值.
\end{enumerate}
\end{problem}

\begin{solution}
\begin{enumerate}
\item 椭圆方程为 \(\dfrac{x^2}{4}+y^2=1\).设直线与椭圆的交点为 \(E=(-x_0,-kx_0)\)、\(F=(x_0,kx_0)\)，其中 \(x_0=\dfrac{2}{\sqrt{1+4k^2}}>0\).直线 \(AB\) 的方程为 \(x+2y=2\)，故
\(D=\left(\dfrac{2}{1+2k},\dfrac{2k}{1+2k}\right)\).

由 \(\overrightarrow{ED}=6\overrightarrow{DF}\)，点 \(D\) 位于 \(E\)、\(F\) 之间，且在横坐标方向上满足 \(x_D-x_E=6(x_F-x_D)\)，所以 \(x_D=\dfrac{x_E+6x_F}{7}=\dfrac{10}{7\sqrt{1+4k^2}}\).于是
\(\dfrac{2}{1+2k}=\dfrac{10}{7\sqrt{1+4k^2}}\)，即 \(7\sqrt{1+4k^2}=5(1+2k)\).两边均为正，平方并整理得 \(24k^2-25k+6=0\)，所以 \(k=\dfrac23\) 或 \(k=\dfrac38\).
\item 仍设 \(E=(-x_0,-y_0)\)、\(F=(x_0,y_0)\)，其中 \(x_0>0\)、\(y_0>0\)，且由椭圆方程知 \(x_0^2+4y_0^2=4\).直线 \(AB\) 为 \(x+2y-2=0\)，且 \(|AB|=\sqrt5\).点 \(E\)、\(F\) 在直线 \(AB\) 两侧，故四边形面积为
\(S_{AEBF}=\dfrac12|AB|(d_E+d_F)=x_0+2y_0\)，其中 \(d_E\)，\(d_F\) 分别为两点到直线 \(AB\) 的距离.

由柯西不等式，\((x_0+2y_0)^2\leqslant2(x_0^2+4y_0^2)=8\)，所以 \(S_{AEBF}\leqslant2\sqrt2\).当 \(x_0=2y_0\) 时取等号，且该条件可由 \(k=\dfrac12\) 实现.因此四边形 \(AEBF\) 面积的最大值为 \(2\sqrt2\).
\end{enumerate}
\end{solution}
